\section{The Joint Bootstrap and the Proof of the Reduction}
\label{sec:bootstrap}

We prove Proposition~\ref{prop:reduction}. The main step is
Proposition~\ref{prop:bootstrap}: for every number $n\ge2$ of rings the clock is
defined up to level $1$, ring $1$ completes its active phase at a time
$T_1\le\Tstar$, where $\Tstar$ does not depend on $n$, and at $T_1$ the material
inside ring $1$ lies in $B_{a/8}$. From this the proof of Proposition~\ref{prop:reduction} obtains a
lower bound, independent of $n$, on a non-radial moment of $\omega^{(n)}(T_1)$,
and applies Corollary~I.\ref{I-cor:defect}.

Throughout this section we assume \textup{(N1)}--\textup{(N3)} and
\textup{(P0)}--\textup{(P6)}, fix $A_1>0$ and an integer $n\ge2$, write
$\omega=\omega^{(n)}$, and use the objects of
Definitions~\ref{def:datum}--\ref{def:clock} with the parameters of
Section~\ref{sec:param}. By Lemma~\ref{lem:nonempty}, (SA) and (Det) hold, so
the results of Section~\ref{sec:relay} apply. We use $\hat a_k$ and $\rhohat_k$
as in Definition~\ref{def:standing}, with the constant $C_0$ of (P4). The
conditions ($\kappa$1)--($\kappa$7) are those of the $\kappa$-list (P5)
(Definition~\ref{def:P5}), and (B1)--(B11) those of the $\sbar$-list (P6)
(Definition~\ref{def:P6}).

Proposition~\ref{prop:bootstrap} is proved by a continuity argument in time
under a hypothesis (BA) bounding every ring by twice what its phase allows; the
lemmas of Section~\ref{sec:relay} return the bounds without the factor $2$. The
clock of Definition~\ref{def:clock} is defined level by level only while the amplitudes $\ksh_k$ lie in
$(0,\delta_0)$, which is one of the conclusions. So we run the argument up to the
time at which the clock would stop, and show at the end that it does not stop.

\begin{definition}
\label{def:stop}
If the clock is defined up to level $1$, put $\Tend=T_1$. Otherwise let
$k_{\mathrm s}$ be the level $k\in\{1,\dots,n\}$ such that the clock is defined up
to level $k+1$ but not up to level $k$ (for $k=n$ the first condition is void),
and put $\Tend=\tsh_{k_{\mathrm s}}$; we then say that the clock \textbf{stops} at
level $k_{\mathrm s}$.
\end{definition}

If the clock is defined up to level $k$, it is defined up to every level
$k'\ge k$; so these levels form a set $\{k_0,\dots,n\}$ or the empty set, and if
$1$ is not in it, there is exactly one level $k_{\mathrm s}$ as in
Definition~\ref{def:stop}: $k_{\mathrm s}=k_0-1$ if the set is
$\{k_0,\dots,n\}$, and $k_{\mathrm s}=n$ if it is empty. By Definition~\ref{def:clock}, for $j<n$, $\tsh_j$
is defined if and only if the clock is defined up to level $j+1$, and $T_j$ if
and only if it is defined up to level $j$. So, if the clock stops, $\tsh_j$ is
defined exactly for $j\ge k_{\mathrm s}$, $T_j$ exactly for $j>k_{\mathrm s}$,
and $S_{k_{\mathrm s}}(\tsh_{k_{\mathrm s}})\not\subset N_{k_{\mathrm s}}$ or
$\ksh_{k_{\mathrm s}}\notin(0,\delta_0)$. In either case $\Tend<\infty$, and by
Lemma~\ref{lem:clock}(1), $[0,\Tend]$ is the union of at most $2n-1$
consecutive \textbf{phase intervals} $\Phase_n=[0,T_n]$, $\Pin_{n-1}$,
$\Phase_{n-1},\dots$, ending with $\Phase_1$ if the clock does not stop, with
$\Pin_{k_{\mathrm s}}=[T_{k_{\mathrm s}+1},\tsh_{k_{\mathrm s}}]$ if it stops at
$k_{\mathrm s}<n$; if $k_{\mathrm s}=n$, then $\Tend=0$.

As in Definition~\ref{def:loc}, ring $j$ is \textbf{quiet} at $t\in[0,\Tend]$ if
$\tsh_j$ is not defined or $t\le\tsh_j$, \textbf{active} at $t$ if $T_j$ is
defined and $\tsh_j\le t\le T_j$, and \textbf{debris} at $t$ if $T_j$ is defined
and $t>T_j$. Every ring has at least one phase at every $t\in[0,\Tend]$: if $T_j$
is not defined, then $j\le k_{\mathrm s}$, and either $\tsh_j$ is not defined or
$j=k_{\mathrm s}$ and $t\le\Tend=\tsh_j$. By Lemma~\ref{lem:clock}(1), if ring
$j$ is quiet or active at $t$, then every ring $i<j$ is quiet at $t$.

\begin{definition}
For $t\in[0,\Tend]$, \textup{(BA)}$(t)$ means that for every $t'\in[0,t]$:
\begin{description}[font=\normalfont]
\item[\textup{(BA1)}] for every ring $j$ that is not debris at $t'$: $S_j(t')\subset\Nflat_j$;
\item[\textup{(BA2)}] for every ring $j<n$ that is quiet at $t'$: the field $h_j$ is
  defined at every time in $[0,t']$, and $\maj_j(A_jt')\le2\Abar\varepsilon_j$;
\item[\textup{(BA3)}] for every ring $k$ that is active at $t'$, with
  $u=A_k(t'-\tsh_k)$: $\|w_k(t')\|_{H^s}\le2C_2\hat a_k(u)(1+\rhohat_k(u))$.
\end{description}
\textup{(BA$^*$)}$(t)$ is the same statement with $\Nflat_j$, $2\Abar$ and $2C_2$
replaced by $\NN_j$, $\Abar$ and $C_2$. \textup{(BA$^\circ$)}$(t)$ is
\textup{(BA)}$(t)$ without, for each level $k$, the instance of \textup{(BA3)}
for ring $k$ at $t'=\tsh_k$.
\end{definition}

In (BA3), $w_k(t')$ is defined, since (BA1) gives
$S_k(t')\subset\Nflat_k\subset N_k$. Since $\NN\subset\operatorname{int}\Nflat$
(Lemma~\ref{lem:annuli}) and $\hat a_k>0$, (BA$^*$)$(t)$ implies (BA)$(t)$ with
room in each inequality. The instance of (BA3) for ring $k$ at $t'=\tsh_k$ is
left out of (BA$^\circ$) because ring $k$ becomes active only at $\tsh_k$, so
this bound cannot be obtained by continuity from bounds at earlier times. The
proof of Proposition~\ref{prop:bootstrap} proves it directly: for $k<n$ in
Step~3, from the hand-off of ring $k$ (Equation~\eqref{eq:boot-switch}), and for
$k=n$ in Step~0 (Equation~\eqref{eq:boot-start}).

\begin{proposition}
\label{prop:bootstrap}
Assume \textup{(N1)}--\textup{(N3)} and \textup{(P0)}--\textup{(P6)}. Then for
every $A_1>0$ and every $n\ge2$:
\begin{enumerate}[label=\textup{(\arabic*)}]
\item $\Tend=T_1$; that is, the clock is defined up to level $1$. Moreover
  $\ksh_1\le\Kup s_1^{m+2}$;
\item the conclusions of Theorem~\ref{thm:active} hold for ring $1$ on
  $\Phase_1$, in ring-$1$ units and active time, with $G=G_1$,
  $\kappa=\ksh_1$, the constant $C_0$ and $\eta_b=\eta_{b,1}$; and
  $|c_1(\hat T_1)|e^{\sigma_1\hat T_1}\ge\frac{15}{16}\delta_0$;
\item $D_2(T_1)\subset B_{a/8}$;
\item $T_1\le\Tstar$, with $\Tstar$ of Lemma~\ref{lem:time}.
\end{enumerate}
\end{proposition}

The next lemma collects the geometric consequences of (BA1) and of lower bounds
on the amplitudes $\ksh_k$: for each ring $j$ that is not debris, the material
inside ring $j$ stays in the small disk $B_{ar_j/8}$ and the rings outside it stay
outside $B_{6r_j}$, so that its external field is defined; and (Loc) holds. Its proof
applies Proposition~\ref{prop:confine} level by level. Step~1 of the proof of
Proposition~\ref{prop:bootstrap} applies the lemma on each phase interval.

\begin{lemma}
\label{lem:geometry}
Let $t\in[0,\Tend]$. Assume that \textup{(BA1)} holds at every $t'\in[0,t]$, and
that $\ksh_k\ge\Klo s_k^{m+2}$ for every level $k<n$ with $\tsh_k<t$. For a level
$k<n$ such that $T_{k+1}$ is defined, put $e_k=T_k$ if $T_k$ is defined and
$e_k=\tsh_k$ otherwise. Then:
\begin{enumerate}[label=\textup{(\alph*)}]
\item for every level $k<n$ such that $T_{k+1}$ is defined and $T_{k+1}\le t$, the
  hypotheses of Proposition~\ref{prop:confine} hold with $t_0=T_{k+1}$ and
  $t_1=\min(t,e_k)$; and for every $t'\in[T_{k+1},\min(t,e_k)]$,
  \[
  D_{k+1}(t')\subset\overline B_{\varrho_k(t')}\subset B_{ar_k/8},\qquad
  |q_{k+1}(t')|\le\Cmom\big(1+A_{k+1}(t'-T_{k+1})\big)^m,
  \]
  where $\varrho_k(t')=\big(2+2\Cstar A_{k+1}(t'-T_{k+1})\big)r_{k+1}$;
\item for every $t'\in[0,t]$ and every level $j$ that is not debris at $t'$, every
  ring $i<j$ is quiet at $t'$ and satisfies
  $S_i(t')\subset\Nflat_i\subset\R^2\setminus B_{6r_j}$, and
  $D_{j+1}(t')\subset B_{ar_j/8}$; hence the field $h_j$ is defined at $t'$, the
  support hypotheses of Lemma~\ref{lem:reduction} hold for ring $j$ in ring-$j$
  units at $t'$, and, if $j<n$, $\eta_j(A_jt')\le1$;
\item \textup{(Loc)} holds at every $t'\in[0,t]$.
\end{enumerate}
\end{lemma}

\begin{proof}
Recall that $\Nflat_i=r_i\Nflat$, with the annulus $\Nflat$ of
Section~\ref{sec:II-prelim}. By (BA1), a ring $i$ that is not debris at $t'\in[0,t]$ has
$S_i(t')\subset\Nflat_i\subset\{\frac58ar_i\le|x|\le\frac74r_i\}$; and
$s_k\le a/32$ for every $k$ by Equation~\eqref{eq:sk-small}.

\emph{Outer rings in \textup{(b)}.} Let $j$ be not debris at $t'\in[0,t]$. Every
ring $i<j$ is quiet at $t'$ (Lemma~\ref{lem:clock}(1)), so
$S_i(t')\subset\{|x|\ge\frac58ar_i\}$, and
$\frac58ar_i\ge\frac58ar_{j-1}=\frac{5a}{8s_{j-1}}r_j\ge20r_j$; hence
$S_i(t')\cap B_{6r_j}=\emptyset$.

The material inside ring $j$ is controlled from the inside out: the confinement
of $D_{j+1}$ after the time $T_{j+1}$, at which the active phase of ring $j+1$
ends, starts from the confinement of $D_{j+2}$ at $T_{j+1}$. So we argue by
downward induction on $j=n,\dots,1$ and show that $D_{j+1}(t')\subset
B_{ar_j/8}$ at every $t'\in[0,t]$ at which ring $j$ is not debris, and, if
$j<n$, that (a) holds for $k=j$. For $j=n$, $D_{n+1}=\emptyset$. Let $j<n$ and
assume the claim for $j+1$.

\emph{\textup{(a)} for $k=j$.} Let $T_{j+1}\le t$ be defined and
$t_1=\min(t,e_j)$; the clock is defined up to level $j+1$. We check the
hypotheses of Proposition~\ref{prop:confine}.
(i): ring $j+1$ is not debris at $T_{j+1}$, so
$S_{j+1}(T_{j+1})\subset\overline B_{7r_{j+1}/4}$, and
$D_{j+2}(T_{j+1})\subset B_{ar_{j+1}/8}$ by induction; so
$D_{j+1}(T_{j+1})\subset\overline B_{2r_{j+1}}$. (ii): for
$t''\in[T_{j+1},t_1]$ and $i\le j$, ring $i$ is not debris at $t''$, since
$t''\le e_j$ and $T_i\ge T_j$ whenever $T_i$ is defined
(Lemma~\ref{lem:clock}(1)); so $S_i(t'')\subset\{|x|\ge\frac a2r_i\}$. (iii) is
(B1) and (B2). (iv): we bound the kick time $A_{j+1}(t_1-T_{j+1})$ of ring $j$
at $t_1$. If $t_1\le\tsh_j$, then $A_{j+1}(t_1-T_{j+1})\le\Tpin$, as
$\tsh_j=T_{j+1}+\Tpin/A_{j+1}$ (Definition~\ref{def:clock}). If $t_1>\tsh_j$,
then $e_j=T_j$, so the clock is defined up to level $j$, and $\tsh_j<t$. So
$\ksh_j\ge\Klo s_j^{m+2}$ by hypothesis, and $\sigma_j\ge\sigma'$ by
Lemma~\ref{lem:detuning}. As $t_1\le T_j=\tsh_j+\hat T_j/A_j$ and
$A_{j+1}=\mu A_j$, $A_{j+1}(t_1-T_{j+1})\le\Tpin+\mu\hat T_j$. It remains to
bound $\hat T_j$, the length of the active phase of ring $j$ in active time: as
$\ksh_j<\delta_0$ (the clock is defined up to level $j$), $\Klo>0$ (Definition~\ref{def:K1}; $c_N>0$ by Lemma~I.\ref{I-lem:uniform}) and
$\sigma_j\ge\sigma'>0$ (SA1),
$\hat T_j=\sigma_j^{-1}\log(\delta_0/\ksh_j)\le\sigma'^{-1}\log\big(\delta_0/(\Klo s_j^{m+2})\big)\le\sigma'^{-1}\big((m+2)\log\frac1{s_j}+\log_+\frac{\delta_0}{\Klo}\big)$.
In both cases, as $\Tpin<\mu(\Tpin+1)$,
\[
A_{j+1}(t_1-T_{j+1})<\mu\Big[\Tpin+1+\frac1{\sigma'}\Big((m+2)\log\frac1{s_j}+\log_+\frac{\delta_0}{\Klo}\Big)\Big],
\]
so $\Lambda=2+2\Cstar A_{j+1}(t_1-T_{j+1})<\Lconf_j$, with $\Lconf_j$ of (B5).
The two inequalities of (iv) increase with $\Lambda$, so they follow from (B5).
Proposition~\ref{prop:confine}(a), (b) give, for $t'\in[T_{j+1},t_1]$,
$D_{j+1}(t')\subset\overline B_{\varrho_j(t')}$ and the bound on $q_{j+1}(t')$,
and $\varrho_j(t')\le\Lambda r_{j+1}<\Lconf_js_jr_j\le\frac a8r_j$ by (B5).

\emph{The claim on $D_{j+1}$.} Let ring $j$ be not debris at $t'\in[0,t]$. If
ring $j+1$ is not debris at $t'$, then $S_{j+1}(t')\subset\overline
B_{7r_{j+1}/4}$, $D_{j+2}(t')\subset B_{ar_{j+1}/8}$ by induction, and
$2r_{j+1}=2s_jr_j\le\frac a{16}r_j$; so $D_{j+1}(t')\subset B_{ar_j/8}$.
Otherwise $T_{j+1}<t'\le t$, and $t'\le e_j$ (if $T_j$ is defined, this is
$t'\le T_j$; otherwise $j=k_{\mathrm s}$ and $t'\le\Tend=\tsh_j$); so (a) for
$k=j$ gives the claim. This completes the induction.

\emph{End of \textup{(b)}.} Let $j$ be not debris at $t'\in[0,t]$. Then
$D_{j+1}(t')\subset B_{ar_j/8}$ and the rings $i<j$ do not meet $B_{6r_j}$, so
$h_j$ is defined at $t'$ (Definition~\ref{def:regions}), and in ring-$j$ units
the inner part is supported in $D_{j+1}(t')/r_j\subset B_{a/8}$ and the outer part in
$\R^2\setminus B_6$, as Lemma~\ref{lem:reduction} requires. If $j<n$, the
hypotheses of Lemma~\ref{lem:field-bound} hold at $t'$, since (B1) holds and
$S_i(t')\subset\Nflat_i\subset N_i$ for $i<j$; so $\eta_j(A_jt')\le
C_\eta(\mu\sbar^2+\sbar^{m-2})\le1$ by (B9).

\emph{\textup{(c)}.} Item (1) of Definition~\ref{def:loc} holds because a ring
$i$ with $t'\le T_i$, or with $T_i$ not defined, is not debris at $t'$, so
$S_i(t')\subset\Nflat_i\subset N_i$. Item (2) concerns a ring $j$ quiet at
$t'$. Its first two conditions, on the rings $i<j$ and on $D_{j+1}(t')$, hold
by (b). Its last clause, the confinement $D_{j+1}(t')\subset\overline
B_{\varrho_j(t')}$ when $T_{j+1}$ is defined and $t'\ge T_{j+1}$, holds by (a)
with $k=j<n$, since a ring $j$ quiet at $t'$ is not debris, so $t'\le e_j$
as above; for $j=n$ it holds as $D_{n+1}=\emptyset$.
\end{proof}

\begin{proof}[Proof of Proposition~\ref{prop:bootstrap}]
Throughout, (SA) and (Det) hold by
Lemma~\ref{lem:nonempty}, and the phases are ordered as in
Lemma~\ref{lem:clock}\textup{(1)}. Step~0 checks the bootstrap assumption at
$t=0$, Steps~1--6 deduce (BA$^*$)$(t)$ from (BA$^\circ$)$(t)$, and Step~7 is the
continuity argument.

\emph{Step 0 (the initial time; it does not use \textup{(BA)}).} At $t=0$, $\tsh_n=0$, $G_n=0$ and, in
ring-$n$ units, $w_n(0)=\Phi^{\kseed}-\Phi$. By ($\kappa$6) at $\kappa=2\kseed$,
$\kseed\le\kL$, so Lemma~\ref{lem:seed} and (P4) give
$S_n(0)=S_n\subset\NNm_n\subset N_n$,
$\ksh_n=|\ellst(0)(\Phi^{\kseed}-\Phi)|\in[\frac12\kseed,2\kseed]$ and
$\|w_n(0)\|_{H^{M-1}}\le\Cseed\kseed\le2\Cseed\ksh_n\le C_0\ksh_n$. By
($\kappa$1) at $\kappa=2\kseed$, $2\kseed<\delta_0$. So $\ksh_n\in(0,\delta_0)$,
the clock is defined up to level $n$, and $k_{\mathrm s}<n$ if it stops.
Recall that $\rhohat_n$ is the transient of Theorem~\ref{thm:active} for
$(C_0,G_n)$; it equals $C_0\jap{u}^pe^{-(\sigma_n-\sigma')u}$ for
$u<t_1(C_0,G_n)$, and $t_1(C_0,G_n)>0$ because $2^{p/2}C_0>1$, so
$\rhohat_n(0)=C_0$. As $s\le M-1$ and $C_2\ge1$,
\begin{equation}
\label{eq:boot-start}
\|w_n(0)\|_{H^s}\le C_0\ksh_n<C_2\ksh_n(1+C_0)=C_2\hat a_n(0)\big(1+\rhohat_n(0)\big).
\end{equation}
For $j<n$, $S_j(0)=S_j\subset\{ar_j\le|x|\le r_j\}\subset r_j\NNm$ and
$\maj_j(0)=0$; $h_j$ is defined at $0$, since
$D_{j+1}(0)\subset\overline B_{5r_{j+1}/4}\subset B_{ar_j/8}$ and the rings $i<j$
lie in $\{|x|\ge ar_i\}\subset\R^2\setminus B_{6r_j}$. Hence (BA$^*$)$(0)$ holds.

Now let $t\in[0,\Tend]$ and assume (BA$^\circ$)$(t)$. We prove the conclusions of
Steps~1--6 for the times in $[0,t]$, treating the phase intervals of $[0,t]$ in
chronological order: Steps~1, 2, 4, 5, 6 on $\Phase_n\cap[0,t]$; then, for
$k=n-1,n-2,\dots$, Steps~1, 2 on $\Pin_k\cap[0,t]$, Step~3 at $\tsh_k$ if
$\tsh_k\le t$, Steps~5, 6 on $\Pin_k\cap[0,t]$, and Steps~1, 2, 4, 5, 6 on
$\Phase_k\cap[0,t]$. Each step uses only (BA$^\circ$)$(t)$ and the steps already
carried out. In particular, Step~3 at $\tsh_k$ uses, besides Step~0, only
Steps~1, 2 on $[0,\tsh_k]$, Step~4 on $[0,T_{k+1}]$ and Step~3 at the levels
$i>k$; and the bound $\ksh_k\ge\Klo s_k^{m+2}$, which Lemma~\ref{lem:quiet-time}
in Step~5 at the time $\tsh_k$ and Step~1 on $\Phase_k$ need, comes from Step~3
at $\tsh_k$.

\emph{Step 1 (geometry).} Let $I$ be the current phase interval and $\bar t\le t$
the right end of $I\cap[0,t]$. Lemma~\ref{lem:geometry} applies with $\bar t$
in place of $t$: (BA1) holds on $[0,\bar t]$, and every level $k<n$ with
$\tsh_k<\bar t$ has $\tsh_k$ at most the left end of $I$, so Step~3 at $\tsh_k$
has already been carried out and has given $\ksh_k\ge\Klo s_k^{m+2}$. Below we use Lemma~\ref{lem:geometry}(a)--(c)
on $[0,\bar t]$ without further mention.

\emph{Step 2 (quiet rings).} Let $t'$ be in the current interval and $j<n$ quiet
at $t'$. If $t'=0$, Equation~\eqref{eq:boot-quiet} below holds at $t''=0$ by
Step~0, as $\maj_j(0)=0$ and $w_j(0)=0$ (ring $j$ starts at $\Phi$ in ring-$j$
units); let $t'>0$ and $\tau=A_jt'$. Suppose first that $\tsh_j$ is defined;
then the clock is defined up to level $j+1$ and $\tau\le\tash_j$. We apply
Lemma~\ref{lem:O} to $j$ and $\tau$: at every $t''\in[0,t']$ every ring $i\le j$
is quiet, so Lemma~\ref{lem:geometry}(b) gives its geometric hypotheses and
$\eta_i(A_it'')\le1$; (BA2) and (B9) give, as $\varepsilon_j=\mu s_j^{m+2}$ and
$s_j\le\sbar$,
$\maj_j\le2\Abar\varepsilon_j\le2\Abar\mu\sbar^{m+2}\le\majst$ on $[0,\tau]$; and
(B7) gives the three smallness conditions on $\sbar$ in Lemma~\ref{lem:O}. So, for $t''\in[0,t']$ and $i<j$,
$\maj_i(A_it'')\le\frac34\Ptil_{i,j}\maj_j(A_jt'')$, the hypotheses of
Theorem~\ref{thm:quiet} hold for ring $i$ on $[0,A_it']$ (for $i=j$ as well), and
$\etout_j\le\theta_j\maj_j$ on $[0,\tau]$. If $\tsh_j$ is not defined, the clock stops,
$j<k_{\mathrm s}$, ring $k_{\mathrm s}$ is quiet on $[0,\Tend]$, and the case just
treated, applied to $k_{\mathrm s}$ at $t'$, shows that Theorem~\ref{thm:quiet}
applies to ring $j$ on $[0,A_jt']$ and
$\maj_j(A_jt')\le\frac34\Ptil_{j,k_{\mathrm s}}\maj_{k_{\mathrm s}}(A_{k_{\mathrm s}}t')$.
In both cases Theorem~\ref{thm:quiet} gives, for $t''\in[0,t']$,
\begin{equation}
\label{eq:boot-quiet}
S_j(t'')\subset\NNm_j,\qquad \|w_j(t'')\|_{H^{s_h}}\le C_Q\maj_j(A_jt'').
\end{equation}

\emph{Step 3 (hand-off of ring $k<n$ at $\tsh_k$, if $\tsh_k\le t$).} This step
starts ring $k$: it shows that the clock is defined up to level $k$, with
$\ksh_k$ between $\Klo s_k^{m+2}$ and $\Kup s_k^{m+2}$, and proves (BA3) for ring
$k$ at $\tsh_k$. Lemma~\ref{lem:kicktrue} bounds the kick amplitude
$|\Ain(\Tpin)|$ that ring $k+1$ and the rings inside it give to ring $k$
(Lemma~\ref{lem:coupling}\textup{(2)}),
Proposition~\ref{prop:relay} bounds the Duhamel majorant $\maj_k(\tash_k)$ from
above, and Lemma~\ref{lem:handoff} combines the two. The clock is
defined up to level $k+1$, Steps~1, 2 have been carried out on $[0,\tsh_k]$ and
Steps~4, 5, 6 on $[0,T_{k+1}]$, and no ring is active on $(T_{k+1},\tsh_k)$; so
Step~4 has been carried out on the whole of $\Phase_{k+1}$ and, if $k+1<n$,
Step~3 for ring $k+1$. By Step~3 for ring $k+1$ and
(B6), or by Step~0 if $k+1=n$,
\begin{equation}
\label{eq:boot-kappa}
0<\ksh_{k+1}\le2\kseed .
\end{equation}
We apply Lemma~\ref{lem:kicktrue} with the constants $C_0$ and $K_1$, and check
its hypotheses in the order in which it lists them. First,
$D_{k+1}\subset B_{ar_k/8}$ on $[0,\tsh_k]$ by
Lemma~\ref{lem:geometry}(b). Second, the conclusions of Theorem~\ref{thm:active} and
Corollary~\ref{cor:active-moment} for ring $k+1$ on $\Phase_{k+1}$ hold by
Step~4 with $\eta_{b,k+1}$, hence with $\eta_b\ge\eta_{b,k+1}$, as the bounds
increase with $\eta_b$. Third, Step~4 for ring $k+1$ verified hypotheses (i)--(iv) of
Lemma~\ref{lem:drift} for ring $k+1$ with $\bar u=\hat T_{k+1}$ and
$\eta=\eta_{b,k+1}\in[0,\eta_b]$, together with $\rT(\Tpin)\le\frac1{64}$ and
$\Cls\eta_{b,k+1}\hat T_{k+1}\le\frac1{32}$. Fourth, if $k+1<n$: on $[0,\tsh_{k+1}]$
$D_{k+2}\subset B_{ar_{k+1}/8}$ and the rings $i\le k$ lie outside $B_{6r_{k+1}}$
by Lemma~\ref{lem:geometry}(b); the hypotheses of Theorem~\ref{thm:quiet} hold for
ring $k+1$ on $[0,\tash_{k+1}]$ by Step~2; $\maj_{k+1}(\tash_{k+1})\le K_1\ksh_{k+1}$ is
Equation~\eqref{eq:handoff} for ring $k+1$; the hypotheses of
Proposition~\ref{prop:confine} with $k+1$ in place of $k$ and $t_1=T_{k+1}$ hold
by Lemma~\ref{lem:geometry}(a) for $k+1$; and
$\ksh_{k+1}\ge\Klo s_{k+1}^{m+2}$ by Step~3 for ring $k+1$. Fifth, the bound
$|q_{k+1}|\le\Cmom(1+\tauK)^m$ for $\tauK\in[0,\Tpin]$ is
Lemma~\ref{lem:geometry}(a) on $\Pin_k$. Finally, conditions (a)--(f) of
Lemma~\ref{lem:kicktrue} are (P1), (P2), ($\kappa$1), ($\kappa$2), ($\kappa$3) at
$\kappa=\ksh_{k+1}$, which Equation~\eqref{eq:boot-kappa} allows, and (B11). So
$|\qs|\in[\frac{15}{16},\frac{17}{16}]2\pi|\pout_*(G_{k+1})|\delta_0$ and
\begin{equation}
\label{eq:boot-kick}
\frac{|\beta_k|\,|\qs|}{10|z|(\mu-1)}\le|\Ain(\Tpin)|\le\frac{3|\beta_k|\,|\qs|\,\mu}{|z|(\mu-1)} .
\end{equation}

Next we apply Proposition~\ref{prop:relay} with the constants $C_0$ and $K_1$,
$j=k$ and $\tau=\tash_k$. By
Lemma~\ref{lem:quiet-time} at the time $\tsh_k$ (its hypothesis
$\Kup\sbar^{m+2}\le\kseed$ is (B6); its hypotheses on the levels
$i\in(k,n)$ hold by Step~3 for these levels, and $\ksh_n\ge\frac12\kseed$ by
Step~0) and Lemma~\ref{lem:theta-quiet},
\begin{equation}
\label{eq:boot-theta}
\theta_k\tash_k\le\theta_k\bar\tau^\sharp_k\le\min\Big\{\frac12,\frac1{4\Cls K_1}\Big\}.
\end{equation}
The hypotheses (H1)--(H4) of Proposition~\ref{prop:relay} hold. (H1): (Loc) holds on
$[0,\tsh_k]$ by Lemma~\ref{lem:geometry}(c). (H2): $\etout_k\le\theta_k\maj_k$ on
$[0,\tash_k]$ by Step~2, and $\theta_k\tash_k\le\frac12$ by
Equation~\eqref{eq:boot-theta}. (H3): if $k+1<n$, the
conclusions of Theorem~\ref{thm:quiet} for ring $k+1$ hold on
$[0,\tash_{k+1}]$ by Step~2. (H4): $\maj_{k+1}(\tash_{k+1})\le K_1\ksh_{k+1}$ by
Step~3 for ring $k+1$; if $k+1=n$, then $\tash_n=0$ and $\maj_n(0)=0$. The
other conditions of (H4), on $\Phase_{k+1}$, hold: $D_{k+2}\subset B_{ar_{k+1}/8}$ by Lemma~\ref{lem:geometry}(b);
(BA3) for ring $k+1$ holds by (BA$^\circ$)$(t)$, and at $\tsh_{k+1}$ by
Equation~\eqref{eq:boot-switch} for ring $k+1$ or by
Equation~\eqref{eq:boot-start}; and
$\|\hin_{k+1}\|_{C^{M_h}(N')}\le C_0\ksh_{k+1}\jap{u}^m$ by Step~4 for ring
$k+1$. Hence, by Proposition~\ref{prop:relay}(ii), Equation~\eqref{eq:relay}
holds at $\tash_k$.

We bound the right side of Equation~\eqref{eq:relay}. Its last term
$\Xi_{k+1}$ collects the quiet phase of ring $k+1$ and the remaining terms of its
active phase. We bound it by ($\kappa$7), after bounding the exponent
$1-\gamma_k$ from below. By
Lemma~\ref{lem:detuning} and its proof, $(1-e_\mu)\sigma\le\sigma_i\le(1+e_\mu)\sigma$
for $i\in\{k,k+1\}$, $e_\mu(\mu+1)\le\frac14(\mu-1)$, $\frac{15}{16}\le1-e_\mu\le1$
and $x_k\le\frac{17}{16}\sigma$. Hence
\begin{gather*}
1-\frac{\sigma_k}{\mu\sigma_{k+1}}\ge1-\frac{1+e_\mu}{\mu(1-e_\mu)}
=\frac{(\mu-1)-e_\mu(\mu+1)}{\mu(1-e_\mu)}\ge\frac34\cdot\frac{\mu-1}\mu,\\
1-\frac{\sigma'}{\sigma_{k+1}}\ge1-\frac{16\sigma'}{15\sigma}=c_\gamma ,
\end{gather*}
and, as $\gamma_k$ is the larger of $\sigma_k/(\mu\sigma_{k+1})$ and
$\sigma'/\sigma_{k+1}$, $1-\gamma_k\ge\min(\frac34\frac{\mu-1}\mu,c_\gamma)\ge\gamma'_\mu$.
As $\ksh_{k+1}<\delta_0$,
$(\ksh_{k+1}/\delta_0)^{1-\gamma_k}\le(\ksh_{k+1}/\delta_0)^{\gamma'_\mu}$, and
($\kappa$7) at $\kappa=\ksh_{k+1}$ gives $\Xi_{k+1}\le\delta_0/(\sigb(\mu-1))$. So
the bracket in Equation~\eqref{eq:relay} with $j=k$ is at most
$(4C_2+1)\delta_0/(\sigb(\mu-1))+C_D(\Tpin)$; with Equation~\eqref{eq:relay}
this gives the first inequality below, and the bound on $x_k$ above with Equation~\eqref{eq:K1} gives the second:
\[
\maj_k(\tash_k)\le2C_M\varepsilon_ke^{x_k\Tpin}\Big[\frac{(4C_2+1)\delta_0}{\sigb(\mu-1)}+C_D(\Tpin)\Big]\le\Abar\,\varepsilon_k .
\]
Next we bound from below the coefficient $\cin_k$ produced by the kick. By
Lemma~\ref{lem:coupling}\textup{(2)}, $\cin_k=e^{\zeta_0\Tpin}\Ain(\Tpin)$ with
$\Re\zeta_0=\sigma_k/\mu=x_k$, so that $|\cin_k|=e^{x_k\Tpin}|\Ain(\Tpin)|$. As
$G_k,G_{k+1}\in[0,\Gone]$, Lemma~I.\ref{I-lem:uniform} gives the first bounds below,
and they give the second with $\beta_k=\beta(G_k)s_k^{m+2}/(2\pi)$ and the bound
on $|\qs|$:
\[
c_N\le|\beta(G_k)||\pout_*(G_{k+1})|\le C_N,\qquad
\tfrac{15}{16}c_Ns_k^{m+2}\delta_0\le|\beta_k||\qs|\le\tfrac{17}{16}C_Ns_k^{m+2}\delta_0 .
\]
With the lower bound in Equation~\eqref{eq:boot-kick}, $\varepsilon_k=\mu s_k^{m+2}$
and $\frac{15}{16}\cdot\frac1{10}=\frac3{32}$,
\[
|\cin_k|\ge e^{x_k\Tpin}\frac{|\beta_k||\qs|}{10|z|(\mu-1)}\ge\frac{3c_N}{32\mu}\cdot\frac{\varepsilon_k\delta_0e^{x_k\Tpin}}{|z|(\mu-1)}
=\frac{c_K\varepsilon_k\delta_0e^{x_k\Tpin}}{|z|(\mu-1)} .
\]
Multiplying by $\frac12K_1$, Equation~\eqref{eq:K1} and the first inequality of
the bound on $\maj_k(\tash_k)$ above give
\[
\frac12K_1|\cin_k|\ge4C_M\varepsilon_ke^{x_k\Tpin}\Big[\frac{(4C_2+1)\delta_0}{\sigb(\mu-1)}+C_D(\Tpin)\Big]\ge2\maj_k(\tash_k).
\]

The hypotheses of Lemma~\ref{lem:handoff} hold: (i) by
Lemma~\ref{lem:geometry}(b) and Step~2; (ii), with the $\theta_k$ of
Lemma~\ref{lem:O}, by Step~2 and Equation~\eqref{eq:boot-theta}; (iii) by the
last inequality; (iv) by $\maj_k(\tash_k)\le\Abar\varepsilon_k\le\Abar\mu\sbar^{m+2}$
and (B11). So Equation~\eqref{eq:handoff} holds at $\tsh_k$. To bound $\ksh_k$
we use its first bound, the identity
$|\cin_k|=e^{x_k\Tpin}|\Ain(\Tpin)|$, Equation~\eqref{eq:boot-kick} and the bounds
on $|\beta_k||\qs|$. As $e^{x_k\Tpin}\ge1$ and $\frac{15}{320}\ge\frac1{22}$,
$\ksh_k\ge\frac12|\Ain(\Tpin)|\ge\frac12\cdot\frac1{10}\cdot\frac{15}{16}\cdot\frac{c_N\delta_0s_k^{m+2}}{|z|(\mu-1)}\ge\Klo s_k^{m+2}$;
and as $x_k\le\frac{17}{16}\sigma$ and
$\frac32\cdot3\cdot\frac{17}{16}=\frac{153}{32}\le5$,
\[
\ksh_k\le\frac32e^{x_k\Tpin}|\Ain(\Tpin)|\le\frac32\cdot3\cdot\frac{17}{16}\cdot\frac{C_N\delta_0\mu\,s_k^{m+2}}{|z|(\mu-1)}e^{17\sigma\Tpin/16}\le\Kup s_k^{m+2}.
\]
Hence
$0<\ksh_k\le\Kup\sbar^{m+2}\le\kseed<\delta_0$ by (B6) and ($\kappa$1), and
$S_k(\tsh_k)\subset\NNm_k\subset N_k$ by Equation~\eqref{eq:boot-quiet}; so the
clock is defined up to level $k$. Finally, as $s\le s_h=M-1$, $C_0\ge\Chand$,
$C_2\ge1$ and $\rhohat_k(0)=C_0$,
\begin{equation}
\label{eq:boot-switch}
\|w_k(\tsh_k)\|_{H^s}\le\Chand\ksh_k\le C_0\ksh_k<C_2\hat a_k(0)\big(1+\rhohat_k(0)\big),
\end{equation}
which is (BA3) for ring $k$ at $t'=\tsh_k$, with a margin. Step~3 uses neither
(BA3) for ring $k$ at $\tsh_k$ nor any time after $\tsh_k$.

\emph{Step 4 (the active phase of ring $k$).} This step verifies the hypotheses
(Init), (Field), (Coef) and (Small) of Theorem~\ref{thm:active} for ring $k$ on
$\Phase_k\cap[0,t]$. Let $t'\in\Phase_k\cap[0,t]$ and
$\bar u=A_k(t'-\tsh_k)\le\hat T_k$. If $t'=\tsh_k$, we need only the two bounds
at the end of this step for $t''=\tsh_k$, and they hold already: $S_k(\tsh_k)\subset\NNm_k\subset\NN_k$ by Step~0
if $k=n$ and by Equation~\eqref{eq:handoff} if $k<n$, and the bound on
$\|w_k(\tsh_k)\|_{H^s}$ is Equation~\eqref{eq:boot-start}, respectively
\eqref{eq:boot-switch}. So let $t'>\tsh_k$; then $\bar u>0$, and $T_k$ is defined
by Step~0 or Step~3. By the invariance of Equation~\eqref{eq:euler} under
translations in time and the scaling of ring-$k$ units, and by
Theorem~\ref{thm:smooth}(5), the solution in ring-$k$ units with time origin
$\tsh_k$ is the solution of Theorem~\ref{thm:smooth} with datum
$A_k^{-1}\omega(\tsh_k,r_k\cdot)$. We apply Lemma~\ref{lem:reduction} and
Theorem~\ref{thm:active} to it, with the parts of the vorticity on $S_k$, on
$D_{k+1}$ and on $S_1\cup\dots\cup S_{k-1}$, with $G=G_k$, $\kappa=\ksh_k$, the
constant $C_0$, $\eta_b=\eta_{b,k}$ and $\delta_0$; then
$w(u)=w_k(\tsh_k+u/A_k)$, $h(u)=h_k(\tash_k+u)$ and $c=c_k$. The support
hypotheses of Lemma~\ref{lem:reduction} on $[0,\bar u]$ hold by
Lemma~\ref{lem:geometry}(b). For Theorem~\ref{thm:active}: $G_k\in[0,\Gone]$,
$\delta_0\le\dH$ by (P1), $\eta_{b,k}\le\eta_b\le\etmax$ by (B8), and (Small)
holds: $\ksh_k\le2\kseed\le\kappa_0(C_0)$ by Step~0 if $k=n$, by Step~3 and (B6) if
$k<n$, and by ($\kappa$4). It remains to check (Init), (Field) and (Coef).

(Init): for $k<n$, Equation~\eqref{eq:handoff} gives
$\|w_k(\tsh_k)\|_{H^{M-1}}\le\Chand\ksh_k\le C_0\ksh_k$ and
$\supp(\Phi+w_k(\tsh_k))\subset\NNm$; for $k=n$ this is Step~0.

(Field), inner part $h_1(u)=\hin_k(\tash_k+u)$. For $k=n$, $h_1=0$, because
$D_{n+1}=\emptyset$ and $G_n=0$. Let $k<n$, $u\in[0,\bar u]$ and
$t''=\tsh_k+u/A_k$. By Lemma~\ref{lem:geometry}(a), $\omega_{\ge k+1}(t'')$ is
supported in $\overline B_R$ with $R=\varrho_k(t'')\le\frac a8r_k$,
$A_{k+1}(t''-T_{k+1})=\Tpin+\mu u$, and $D_{k+1}(t'')\subset B_{ar_k/8}$, so
$h_1(u)$ is the field $E_{[k]}$ of Lemma~\ref{lem:inner-field} at $t''$; and
$\|\omega_{\ge k+1}(t'')\|_{L^1}\le\Mvort A_{k+1}r_{k+1}^2$ by
Lemma~\ref{lem:speed}. Lemma~\ref{lem:inner-field}, the inequalities
$2+2\Cstar\Tpin+2\Cstar\mu u\le(2+2\Cstar\Tpin)(1+\Cstar\mu u)$ and
$1+\Cstar\mu u\le(1+\Cstar\mu)\max(1,u)\le(1+2\Cstar)\jap{u}$, the definition of
$C_F$ in (P3), and $\ksh_k\ge\Klo s_k^{m+2}$ (Step~3) give
\begin{equation}
\label{eq:boot-inner}
\begin{aligned}
\|h_1(u)\|_{C^{M_h}(N')}&\le\Cin\mu s_k^2(R/r_k)^m=\Cin\mu s_k^{m+2}\big(2+2\Cstar(\Tpin+\mu u)\big)^m\\
&\le\Cin\mu(2+2\Cstar\Tpin)^m(1+2\Cstar)^ms_k^{m+2}\jap{u}^m\\
&=C_F\Klo s_k^{m+2}\jap{u}^m\le C_F\ksh_k\jap{u}^m .
\end{aligned}
\end{equation}
As $M\le M_h$ and $C_F\le C_0$, this is the first bound of (Field).

(Field), outer part $h_2(u)=\hout_k(\tash_k+u)$. We apply
Lemma~\ref{lem:Oprime} with the constants $C_0$, $K_1$, $C_F$ at the time $t'$.
The clock is defined up to level $k$, and (Loc) holds on $[0,t']$ by
Lemma~\ref{lem:geometry}(c). If $k\ge2$, ring $k-1$ is quiet on $[0,t']$, so
Step~2 gives the hypotheses and conclusions of Lemma~\ref{lem:O} for $j=k-1$ and
$\tau=A_{k-1}t'$; and $\theta_{k-1}A_{k-1}t'\le\theta_{k-1}\tash_{k-1}\le\frac12$
(as ring $k-1$ is quiet at $t'$) by Lemma~\ref{lem:quiet-time} at the time
$\tsh_{k-1}$ and Lemma~\ref{lem:theta-quiet}; the hypotheses of
Lemma~\ref{lem:quiet-time} besides (B6) concern the levels $i\ge k$, whose amplitudes are
bounded below by Step~3 at the times $\tsh_i\le\tsh_k$ and, for $i=n$, by
Step~0. (BA3) for ring $k$ holds on $(\tsh_k,t']$ by (BA$^\circ$)$(t)$ and at
$\tsh_k$ by Equation~\eqref{eq:boot-switch} or \eqref{eq:boot-start}; the inner
bound is Equation~\eqref{eq:boot-inner}. If $k<n$, the conclusions of
Theorem~\ref{thm:quiet} for ring $k$ on $[0,\tash_k]$ hold by Step~2,
$\maj_k(\tash_k)\le K_1\ksh_k$ by Step~3, and $D_{k+1}\subset B_{ar_k/8}$ on
$[\tsh_k,t']$ by Lemma~\ref{lem:geometry}(a). So
$\|h_2(u)\|_{C^{M_h}(N')}\le\eta_{b,k}\hat a_k(u)$ on $[0,\bar u]$, which is the
second bound of (Field).

(Coef): we apply Lemma~\ref{lem:drift} to ring $k$ with this $\bar u$ and
$\eta=\eta_{b,k}$. Hypotheses (i) and (ii) concern only $k<n$. (i) is
Lemma~\ref{lem:geometry}(a) on $W_k$, since the kick time is $\tauK=\Tpin+\mu u$.
(ii), with $q'=\qs$ of Step~3: the hypotheses of Lemma~\ref{lem:coupling}(2) hold
by Lemma~\ref{lem:geometry}(b), so Equation~\eqref{eq:handoff} and
Lemma~\ref{lem:coupling}(2) give
$\ksh_k\ge\frac12|\cin_k|=\frac12e^{x_k\Tpin}|\Ain(\Tpin)|$; Step~3 gives
the lower bound on $|\Ain(\Tpin)|$. By Step~3 and Lemma~I.\ref{I-lem:uniform},
$|\qs|\ge\frac{15}{16}2\pi|\pout_*(G_{k+1})|\delta_0\ge\frac{15}{16}\cout\delta_0>0$,
so $\qs\ne0$, and (P2) gives
$\mu-1\le c'\delta_0\le|\qs|/(10|z|\Cmom\Ip{m}{\sigma/4})$. (iii) is the outer
bound of (Field), the $C^0$ norm being at most the $C^{M_h}$ norm, and (iv) is
Lemma~\ref{lem:geometry}(b). Moreover $\rT(\Tpin)\le\frac1{64}$ by (P3), and
$\Cls\eta_{b,k}\hat T_k\le\frac1{32}$: for $k=1$, $\eta_{b,1}=0$; for $2\le k<n$,
$\hat T_k\le\sigma'^{-1}\big((m+2)\log\frac1{s_k}+\log_+\frac{\delta_0}{\Klo}\big)$,
as $\sigma_k\ge\sigma'$ and $\ksh_k\ge\Klo s_k^{m+2}$, and (B8) applies; for
$k=n$, $\hat T_n\le\sigma'^{-1}\log(2\delta_0/\kseed)$ by Step~0 and
$\eta_{b,n}\le\eta_b$, and (B8) applies. So Lemma~\ref{lem:drift} gives (Coef) on
$[0,\bar u]$ and, if $\bar u=\hat T_k$,
$|c_k(\hat T_k)|e^{\sigma_k\hat T_k}\in[\frac{15}{16}\delta_0,\frac{17}{16}\delta_0]$.

Theorem~\ref{thm:active} and Corollary~\ref{cor:active-moment} therefore hold on
$[0,\bar u]$: in particular $S_k(t'')\subset\NN_k$ and
$\|w_k(t'')\|_{H^s}\le C_2\hat a_k(u)(1+\rhohat_k(u))$ for
$t''=\tsh_k+u/A_k\in[\tsh_k,t']$.

\emph{Step 5 (quiet size).} Let $t'$ be in the current interval, $j<n$ quiet at
$t'$ and $\tau=A_jt'$. We show $\maj_j(\tau)\le\Abar\varepsilon_j$; by
Lemma~\ref{lem:geometry}(b), $h_j$ is defined on $[0,t']$. If $t'=0$, then
$\maj_j(0)=0$ (Step~0); let $t'>0$. Suppose first that
$\tsh_j$ is defined, so that $\tau\le\tash_j$, and apply
Proposition~\ref{prop:relay} with the constants $C_0$ and $K_1$ to $j$ and
$\tau$. (Loc) holds on $[0,t']$ by
Lemma~\ref{lem:geometry}(c), and $\etout_j\le\theta_j\maj_j$ on $[0,\tau]$ by
Step~2. Lemma~\ref{lem:quiet-time} at the time $t'$, whose hypotheses on the
levels $i\in(j,n)$ with $\tsh_i\le t'$ hold by Step~3 at the times $\tsh_i\le t'$
and $\ksh_n\ge\frac12\kseed$ by Step~0, gives $\tau\le\bar\tau^\sharp_j$, and
Lemma~\ref{lem:theta-quiet} gives $\theta_j\bar\tau^\sharp_j\le\frac12$; so
$\theta_j\tau\le\frac12$. If $j+1<n$, the conclusions of Theorem~\ref{thm:quiet}
for ring $j+1$ on $[0,A_{j+1}\min(t',\tsh_{j+1})]$ hold by Step~2. If
$t'>\tsh_{j+1}$, the remaining hypotheses hold exactly as in Step~3 with $k=j$.

If $\tau\le A_j\tsh_{j+1}$, then $j+1<n$, as $\tsh_n=0<t'$, and ring $j+1$ is
quiet at $\tau/A_j$; so Proposition~\ref{prop:relay}(i), (BA2) for ring $j+1$ at the time $\tau/A_j$ and
(B10) give, as $\varepsilon_{j+1}\le\mu\sbar^{m+2}$,
\[
\maj_j(\tau)\le2C_M\Big(\frac{2C_Q}{\sigb(\mu-1)}+1\Big)\varepsilon_j\cdot2\Abar\varepsilon_{j+1}
\le\Abar\varepsilon_j .
\]
If $\tau>A_j\tsh_{j+1}$, Proposition~\ref{prop:relay}(ii) or (iii) bounds
$\maj_j(\tau)$ by the right side of Equation~\eqref{eq:relay}. As in Step~3 with
$k=j$ ($1-\gamma_j\ge\gamma'_\mu$ and $\ksh_{j+1}<\delta_0$), $\Xi_{j+1}\le\delta_0/(\sigb(\mu-1))$ by ($\kappa$7),
and $x_j\le\frac{17}{16}\sigma$ by Lemma~\ref{lem:detuning}; so that right side
is at most
\[
2C_M\varepsilon_je^{17\sigma\Tpin/16}\Big[\frac{(4C_2+1)\delta_0}{\sigb(\mu-1)}+C_D(\Tpin)\Big]=\Abar\varepsilon_j .
\]
Suppose now that $\tsh_j$ is not defined. Then $j<k_{\mathrm s}$, and by Step~2
and the case just treated for ring $k_{\mathrm s}$, which is quiet at $t'$,
\[
\maj_j(\tau)\le\tfrac34\Ptil_{j,k_{\mathrm s}}\maj_{k_{\mathrm s}}(A_{k_{\mathrm s}}t')
\le\tfrac34\Abar\,\KO\varepsilon_j\,\varepsilon_{k_{\mathrm s}}\prod_{p=j+1}^{k_{\mathrm s}-1}\KO\varepsilon_p\le\tfrac38\Abar\varepsilon_j,
\]
because $\KO\varepsilon_p\le\KO\mu\sbar^{m+2}\le\frac12$ for every $p$ by (B7).
So (BA2) holds with $\Abar$ in place of $2\Abar$.

\emph{Step 6 (localization).} Let $j$ be not debris at $t'\in[0,t]$. If $j$ is
quiet at $t'$, then $S_j(t')\subset\NNm_j\subset\NN_j$ by
Equation~\eqref{eq:boot-quiet} (for $j=n$, quiet only at $t'=0$, by Step~0); if
$j$ is active at $t'$, then $S_j(t')\subset\NN_j$ by Step~4. So (BA1) holds with
$\NN_j$ in place of $\Nflat_j$. With Step~5, and with Step~4 and
Equations~\eqref{eq:boot-switch}, \eqref{eq:boot-start} for (BA3): if
(BA$^\circ$)$(t)$ holds, then (BA)$(t)$ and (BA$^*$)$(t)$ hold.

\emph{Step 7 (continuity, and the clock does not stop).} Let $\mathcal B$ be the
set of $t\in[0,\Tend]$ such that (BA$^\circ$)$(t)$ holds. It contains $0$ by
Step~0, and with every $t$ it contains $[0,t]$. Let $t_*=\sup\mathcal B$. We use
three continuity properties. The flow $\flow$ is continuous on
$[0,\Tend]\times\R^2$ (Theorem~\ref{thm:smooth}(2),(4)), so for a compact set $S$
and an open set $U$ with $\flow_{t_0}(S)\subset U$ we have $\flow_{t'}(S)\subset U$
for $t'$ near $t_0$. On an interval on which $S_k(t')\subset N_k$, the map
$t'\mapsto w_k(t')\in H^s$ is continuous, because
$\omega\mathbf 1_{S_k(t')}=f^{(n)}_k\circ\flow_{t'}^{-1}$
(Equation~\eqref{eq:omega-ge}) is smooth in $(t',x)$ with support in a fixed
compact set. On an interval on which $h_j$ is defined, $\tau\mapsto
h_j(\tau)\in C^{M_h}(N')$ is continuous (Lemma~\ref{lem:reduction}(1)), so
$\eta_j$ and $\maj_j$ are continuous.

\emph{$t_*\in\mathcal B$.} The conditions of (BA$^\circ$)$(t_*)$ at times
$t'<t_*$ hold, because $\mathcal B\supset[0,t')$; so let $t_*>0$ and consider
$t'=t_*$. A ring $j$ that is not debris at $t_*$ is not debris on $[0,t_*]$, and
$S_j(t'')\subset\Nflat_j$ for $t''<t_*$; as $\Nflat_j$ is closed,
$S_j(t_*)\subset\Nflat_j$. So (BA1) holds on $[0,t_*]$. Every level $k<n$ with
$\tsh_k<t_*$ has $\tsh_k\in\mathcal B$, so Step~3 gives $\ksh_k\ge\Klo s_k^{m+2}$,
and Lemma~\ref{lem:geometry} applies with $t=t_*$; so, for a ring $j<n$ quiet at
$t_*$, $h_j$ is defined and $\maj_j$ is continuous on $[0,t_*]$, and
$\maj_j(A_jt_*)\le2\Abar\varepsilon_j$ follows from the same bound at earlier
times. Let $k$ be active at $t_*$ with $\tsh_k<t_*$. Then $w_k$ is continuous on
$[\tsh_k,t_*]$, and the right side $R_k(u)=2C_2\hat a_k(u)(1+\rhohat_k(u))$ of
(BA3) is continuous except at $u=t_1(C_0,G_k)$, where it is right-continuous
and jumps upward: the transient $\rhohat_k$ of Theorem~\ref{thm:active} takes the value
$1$ there, while its left limit there is the value of
$C_0\jap{u}^pe^{-(\sigma_k-\sigma')u}$ at $u=t_1(C_0,G_k)$, which is at most
$2^{-p/2}$ by the definition of $t_1(C_0,G_k)$.
So $R_k(u_*)\ge\limsup_{u\uparrow u_*}R_k(u)$ at $u_*=A_k(t_*-\tsh_k)$, and (BA3)
at $t_*$ follows from (BA3) at earlier times. Hence $t_*\in\mathcal B$.

\emph{$t_*=\Tend$.} Suppose $t_*<\Tend$. By Steps~1--6, (BA$^*$)$(t_*)$ holds.
Choose $\Delta>0$ such that $t_*+\Delta\le\Tend$ and no $\tsh_i$ or $T_i$ lies
in $(t_*,t_*+\Delta]$. Then a ring that is not debris (respectively quiet,
active) at a time in $(t_*,t_*+\Delta]$ is not debris (respectively quiet, active
with $\tsh_k\le t_*$) at $t_*$. Each $S_j(t_*)$ with $j$ not debris at $t_*$ is a
compact subset of $\NN_j\subset\operatorname{int}\Nflat_j$, so after decreasing
$\Delta$, (BA1) holds on $[0,t_*+\Delta]$. No $\tsh_k$ lies in
$(t_*,t_*+\Delta]$, so Lemma~\ref{lem:geometry} applies with $t=t_*+\Delta$, and
$h_j$ is defined on $[0,t_*+\Delta]$ for every ring $j<n$ quiet at $t_*+\Delta$;
as $\maj_j(A_jt_*)\le\Abar\varepsilon_j<2\Abar\varepsilon_j$ and $\maj_j$ is
continuous, after decreasing $\Delta$, (BA2) holds on $[0,t_*+\Delta]$. For a
ring $k$ active at times in $(t_*,t_*+\Delta]$,
$\|w_k(t_*)\|_{H^s}\le C_2\hat a_k(1+\rhohat_k)<R_k$ at $u_*$, because
$\hat a_k>0$; as $w_k$ is continuous and $R_k$ is right-continuous, after
decreasing $\Delta$, (BA3) holds on $(t_*,t_*+\Delta]$. So
$t_*+\Delta\in\mathcal B$, contradicting $t_*=\sup\mathcal B$. Hence $t_*=\Tend$
and $\mathcal B=[0,\Tend]$, and by
Steps~1--6, (BA)$(t)$ and (BA$^*$)$(t)$ hold for every $t\in[0,\Tend]$.

If the clock were not defined up to level $1$, then $k_{\mathrm s}<n$ by Step~0
and $\Tend=\tsh_{k_{\mathrm s}}\in\mathcal B$, so Step~3 for $k=k_{\mathrm s}$
would show that the clock is defined up to level $k_{\mathrm s}$, contrary to
Definition~\ref{def:stop}. Hence $\Tend=T_1$.

We read off the items. (1): $\Tend=T_1$; the bound on $\ksh_1$ is Step~3 for
$k=1$. (2): Step~4 at $t=T_1$ gives the conclusions of
Theorem~\ref{thm:active} on all of $\Phase_1$, and Lemma~\ref{lem:drift} with
$\bar u=\hat T_1$ gives the bound on $|c_1(\hat T_1)|e^{\sigma_1\hat T_1}$.
(3): by (BA)$(T_1)$ and Step~3, Lemma~\ref{lem:geometry} applies with $t=T_1$,
and its part (b) with $t'=T_1$ and $j=1$, ring $1$ being active at $T_1$, gives
$D_2(T_1)\subset B_{ar_1/8}=B_{a/8}$, as $r_1=1$. (4): the clock is defined up to
level $1$, $\ksh_k\ge\Klo s_k^{m+2}$ for $k<n$ (Step~3) and $\ksh_n\ge\frac12\kseed$ (Step~0); so
Lemma~\ref{lem:time} gives $T_1\le\Tstar$.
\end{proof}

\begin{proof}[Proof of Proposition~\ref{prop:reduction}]
Assume \textup{(N1)}--\textup{(N3)}. By Lemma~\ref{lem:nonempty}, we may choose
parameters satisfying (P0)--(P6). This fixes $\mu\in(1,\mu_{\max}]$ with
$\mu_{\max}\le\frac{17}{16}$ (by (P2)) and $\sbar\in(0,a/16]$ (by (P6) and
(B1)), and none of (P0)--(P6) involves $A_1$. Let $A_1>0$, let $f$ be the
function of Definition~\ref{def:datum} with these parameters, which is the
function $f$ of Equation~\eqref{eq:datum}, and let $\Tstar$ be the number of
Lemma~\ref{lem:time}. By Proposition~\ref{prop:integrability}, $f$ is radial
and nonnegative, $f\in (L^1\cap L^q)(\R^2)$ for every $q\in[1,\infty)$, and
$f\notin L^\infty(\R^2)$. The bracket in the definition of $\Tstar$ is a finite
sum of the term $1$ and of nonnegative terms, because $\Tz>0$, $\mu>1$,
$0<\sbar\le a/16<1$ and $\kseed<\delta_0$ by ($\kappa$1) at $\kappa=2\kseed$;
so $\Tstar\in[1/A_1,\infty)$, and in particular $\Tstar\in(0,\infty)$.

Fix $\chiout\in C_c^\infty((a/4,3))$ with $0\le\chiout\le1$ and $\chiout=1$ on
$[a/2,2]$, and put $\Psione(x)=\chiout(|x|)\bar x^{\,m}$; it does not depend on
$n$. The cutoff $\chiout$ is chosen so that, at the time $T_1$ below, it equals
$1$ on the support of ring $1$ and vanishes on the material inside ring $1$; so
the moment $\int\omega^{(n)}(T_1)\Psione$ sees only ring $1$.

For $n\ge2$ let $f^{(n)}$ and $\omega^{(n)}$ be as in Definition~\ref{def:approx},
with the seed size $\kseed$ of (P5), which lies in $(0,\kL]$ by ($\kappa$6) at
$\kappa=2\kseed$. By Lemma~\ref{lem:adm}, $(f^{(n)},\omega^{(n)})_{n\ge2}$ is a
$2$-admissible sequence for $f$ on $[0,\infty)$ (indexed from $n=2$), and
$\omega^{(n)}(t)$ is $3$-fold rotationally symmetric for every $n$ and $t$.
For $n\ge2$ put $T^{(n)}=T_1$, which is defined by
Proposition~\ref{prop:bootstrap}(1) and lies in $[0,\Tstar]$ by
Proposition~\ref{prop:bootstrap}(4). (SA) and (Det) hold by
Lemma~\ref{lem:nonempty}, and the clock is defined up to level $1$. By
Proposition~\ref{prop:bootstrap}(2), the conclusions of
Theorem~\ref{thm:active}(a)--(b) hold for ring $1$ on $\Phase_1$ with $G=G_1$,
$\kappa=\ksh_1$, the constant $C_0$ and $\eta_b=\eta_{b,1}=0$
(Definition~\ref{def:P6}). Ring $1$ has
$r_1=1$, so ring-$1$ units are the original lengths; at $t=T_1$ its active time
is $\hat T_1$, so $w(\hat T_1)=w_1(T_1)$ in Theorem~\ref{thm:active}. The
vorticity $\omega^{(n)}(T_1)$ is supported in $D_1(T_1)$, the disjoint union of
$S_1(T_1)$ and $D_2(T_1)$. By Proposition~\ref{prop:bootstrap}(3),
$D_2(T_1)\subset B_{a/8}$, so the part on $D_2(T_1)$ does not
contribute, as $\chiout$ vanishes on $[0,a/4]$. On $S_1(T_1)$,
$\omega^{(n)}(T_1)=A_1(\Phi+w_1(T_1))$ (Definition~\ref{def:regions}), which is
supported in $\NN\subset\{a/2\le|x|\le2\}$ by Theorem~\ref{thm:active}(a), where
$\chiout(|x|)=1$. As $\Phi$ is radial, $\pout_m[\Phi]=0$, with the
multipole moment $\pout_m$ of Definition~I.\ref{I-def:multipole}, so
$\int\omega^{(n)}(T_1)\Psione=A_1\int(\Phi+w_1(T_1))\,\bar y^{\,m}\dd y=2\pi A_1\,\pout_m[w_1(T_1)]$.

Write $c=c_1(\hat T_1)$, $\lambda=\lambda_1$ and $\phist=\phist(G_1)$, the
unstable eigenfunction of Definition~I.\ref{I-def:eigendata}. We bound
$|\pout_m[w_1(T_1)]|$ from below by comparing $w_1(T_1)$ with
$2\Re(ce^{\lambda\hat T_1}\phist)$, as in Theorem~\ref{thm:active}\textup{(b)}. As
$\pout_m[\phist]=\pout_*(G_1)$ and $\pout_m[\overline{\phist}]=0$ (its integrand
carries the factor $e^{-6i\theta}$),
$\pout_m[2\Re(ce^{\lambda\hat T_1}\phist)]=ce^{\lambda\hat T_1}\pout_*(G_1)$. The
function $w_1(T_1)-2\Re(ce^{\lambda\hat T_1}\phist)$ is supported in $N$, because
$\Phi$, $\Phi+w_1(T_1)$ and $\phist\in X^0$ are. (P1) gives
$\Cp C_2\delta_0\le\cout/(16\pi)$, and ($\kappa$5) with $G=G_1\in[0,\Gone]$
(Lemma~\ref{lem:inner-circ} and (SA3)) and
$\kappa=\ksh_1$, which lies in $(0,\kseed]$ by
Proposition~\ref{prop:bootstrap}(1) and (B6) (it is positive since the clock is
defined up to level $1$, Definition~\ref{def:clock}), gives
$\Cp C_2(\rhohat_1(\hat T_1)+\ksh_1{}^{1/2})\le\cout/(16\pi)$. As
$\hat a_1(\hat T_1)=\delta_0$, Lemma~I.\ref{I-lem:uniform},
Theorem~\ref{thm:active}(b) with $\eta_b=0$ and these two inequalities give
\begin{align*}
\big|\pout_m[w_1(T_1)]-ce^{\lambda\hat T_1}\pout_*(G_1)\big|
&\le\Cp\big\|w_1(T_1)-2\Re(ce^{\lambda\hat T_1}\phist)\big\|_{H^s}\\
&\le\Cp C_2\delta_0\big[\rhohat_1(\hat T_1)+\delta_0+\ksh_1{}^{1/2}\big]
\le\frac{\cout\delta_0}{8\pi}.
\end{align*}
As $G_1\in[0,\Gone]$, $|\pout_*(G_1)|\ge\cout/(2\pi)$ (Lemma~I.\ref{I-lem:uniform}), and
$|ce^{\lambda\hat T_1}|=|c_1(\hat T_1)|e^{\sigma_1\hat T_1}\ge\frac{15}{16}\delta_0$
by Proposition~\ref{prop:bootstrap}(2). Hence
$|\pout_m[w_1(T_1)]|\ge\frac{\cout}{2\pi}\cdot\frac{15}{16}\delta_0-\frac{\cout\delta_0}{8\pi}
=\frac{11}{16}\cdot\frac{\cout\delta_0}{2\pi}\ge\frac{\cout\delta_0}{4\pi}$, and
multiplying by $2\pi A_1$ gives
$|\int\omega^{(n)}(T^{(n)})\Psione|\ge c_0:=\frac12A_1\cout\delta_0>0$. The number $c_0$
involves only $A_1$, $\cout$ and $\delta_0$, and $\Tstar$ only $A_1$, $\mu$,
$\sbar$, $\sigma'$, $\Tz$, $\delta_0$ and $\kseed$; none of them depends on $n$.
Moreover $\Psione$ has the form required in
Corollary~I.\ref{I-cor:defect}, because $\chiout\in C_c^\infty((a/4,3))$ is real.
Corollary~I.\ref{I-cor:defect} with $\bar T=\Tstar$ gives a renormalized weak
solution $\omega_\infty$ on $[0,\infty)$ with datum $f$, in the sense of
Definition~I.\ref{I-def:weak} with $q=2$, such that
$\omega_\infty\in C([0,\infty);L^p(\R^2))$ for every $p\in[1,\infty)$ and, for
every $t\ge0$, $\omega_\infty(t)$ is equimeasurable with $f$ and $3$-fold
rotationally symmetric, and a time $t_0\in[0,\Tstar)$ such that $\omega_\infty$ and $f$ are
distinct elements of $X_q(T')$ for every $T'>t_0$ and every $q\in(4/3,\infty)$.

Let $T>t_0$ and $q\in(4/3,\infty)$; then $T>0$. By
Lemma~I.\ref{I-lem:radial-stationary}, the stationary function $t\mapsto f$ is a
renormalized weak solution in $X_q(T)$ with datum $f$. Also $\omega_\infty\in
X_q(T)$, since, by equimeasurability, $\|\omega_\infty(t)\|_{L^r}=\|f\|_{L^r}$ for every
$r\in[1,\infty)$ and $\omega_\infty\in C([0,T];L^1\cap L^q)$. By
Definition~I.\ref{I-def:weak}, the restriction of $\omega_\infty$ to $[0,T)$
satisfies Equations~(I.\ref{I-eq:weak}) and (I.\ref{I-eq:renorm}) for every test function
$\varphi\in\DT$ and every renormalization $\Ren$; these identities involve only
$\omega_\infty$, its velocity $\KBS\ast\omega_\infty$ and the datum $f$, and not
the exponent of the space. So $\omega_\infty$ is a renormalized weak solution in
$X_q(T)$ with datum $f$, and, as $T>t_0$, the two solutions are distinct elements
of $X_q(T)$. This holds in particular for every $T>\Tstar$, since $t_0<\Tstar$.
So $f$, $\Tstar$, the stationary solution, $\omega_\infty$ and $t_0$ have every
property stated in Theorem~\ref{thm:main}, for every $A_1>0$ with the same $\mu$
and $\sbar$, which were fixed before $A_1$.
\end{proof}
